\chapter{Your own neighbourhood}\label{ch:practice}

\section{Reading responsibly}
\input{generated/nb_urban_form_12}

\section{The same measures in your browser}
The \href{https://ahmadammar.github.io/urban-form-explorer/}{Urban form explorer} computes the measures of this guide for any area you choose, with no installation. It fetches OpenStreetMap data from the same public Overpass servers, works in a local metric projection in the browser, and reports streets (density, junctions, dead ends, circuity, route directness, orientation order with a rose diagram), buildings (count, footprints, floors where tagged, GSI, FSI, open space ratio), access (walking distance along streets to the nearest place of each everyday kind) and composite indices for walkability, cyclability and public transport, always shown with their parts. Each figure links to its formula on the methods page, and each analysis reports how much of the area carries the tags it relies on. Use the notebook to understand a measure, and the explorer to apply it widely; when the two disagree, the notebook shows you why.

\section{What you have learned}
\input{generated/nb_urban_form_13}
